\documentclass[UTF8,fontset=fandol,11pt]{ctexart}
\usepackage[a4paper,margin=25mm]{geometry}
\usepackage{amsmath,amssymb,amsthm,mathtools}
\usepackage{xcolor}
\usepackage{booktabs,array,enumitem,fancyhdr,hyperref}
\hypersetup{colorlinks=true,linkcolor=blue!50!black,urlcolor=blue!50!black}
\setlength{\emergencystretch}{2em}
\setlength{\parskip}{0.3em}
\linespread{1.16}
\allowdisplaybreaks
\newtheorem{theorem}{定理}[section]
\newtheorem{proposition}[theorem]{命题}
\newtheorem{lemma}[theorem]{引理}
\theoremstyle{definition}
\newtheorem{definition}[theorem]{定义}
\theoremstyle{remark}
\newtheorem{remark}[theorem]{说明}
\newcommand{\dd}{\,\mathrm d}
\newcommand{\ii}{\mathrm i}
\newcommand{\Rea}{\operatorname{Re}}
\newcommand{\Ima}{\operatorname{Im}}
\newcommand{\Li}{\operatorname{Li}}
\newcommand{\supp}{\operatorname{supp}}
\pagestyle{fancy}
\fancyhf{}
\fancyhead[L]{\small 学习笔记}
\fancyhead[R]{\small 初稿}
\fancyfoot[C]{\thepage}
\setlength{\headheight}{15pt}
\author{}
\date{2026年10月22日初稿}
\makeatletter
\renewcommand{\maketitle}{%
  \thispagestyle{plain}%
  \begin{center}
    {\LARGE\bfseries \@title\par}
    \vspace{0.6em}
    {\small \@date\par}
  \end{center}
  \vspace{0.6em}}
\makeatother
\newcommand{\draftnotice}{\noindent{\small
本文为供逐篇审阅的学习笔记初稿，不是新增研究成果。
文中区分直接证明与引用结果；正文不依赖本文。
本轮已作公式、条件及排版初步核查，尚待作者逐句审定。}\par\medskip}

\title{Li 判据：从零点位置到非负性}
\begin{document}
\maketitle
\draftnotice
\section{把一个熟悉的问题改写为另一个问题}
RH 约束全部非平凡零点的实部。Li 判据将它等价地改写为一列实数的非负性。
本篇的任务不是重述判据的名称，而是解释这列数怎样定义、RH 为什么使其非负，
以及有限计算怎样受到尾项和未知零点的限制。
完整的反向等价定理引用 Li 的原论文；本篇不把尚未重证的反向部分伪装成一个显然推论。

\section{先定义求和的含义}
记 $\rho=\beta+\ii\gamma$ 为 $\zeta$ 的非平凡零点，按重数计。
已知 $0<\beta<1$、$\gamma\ne0$、零点共轭对称，且 $N(T)=O(T\log(T+2))$。
对 $n\ge1$，定义
\[
 \lambda_n=\lim_{T\to\infty}
 \sum_{|\Ima\rho|\le T}
 \left[1-\left(1-\frac1\rho\right)^n\right].
\]
这一高度对称的约定是定义的一部分，不表示未配对级数绝对收敛。

对固定 $n$，二项式展开给出
\[
 1-(1-1/\rho)^n=\frac n\rho+O_n(|\rho|^{-2}).
\]
把 $\rho$ 与 $\bar\rho$ 配对，首项为
$2n\beta/|\rho|^2$，所以整个配对项为 $O_n(\gamma^{-2})$。
由零点计数估计，$\sum_{\gamma>0}\gamma^{-2}<\infty$。
故按共轭对配对后绝对收敛，极限存在且为实数。
这证明了所写定义有意义，而非只凭“共轭配对”四字略过收敛问题。

\section{临界线为何产生非负项}
直接计算
\[
 \left|1-\frac1\rho\right|^2
 =\frac{(\beta-1)^2+\gamma^2}{\beta^2+\gamma^2}
 =1+\frac{1-2\beta}{|\rho|^2}.
\]
因此其模为 $1$ 当且仅当 $\beta=1/2$。
在 RH 下可写 $1-1/\rho=e^{\ii\theta_\rho}$，共轭对贡献变为
\[
 [1-e^{\ii n\theta_\rho}]+[1-e^{-\ii n\theta_\rho}]
 =2(1-\cos(n\theta_\rho))\ge0.
\]
上一节已证明该配对级数收敛，所以 $\lambda_n\ge0$。
这里没有把一个未定义的“无穷个非负项之和”当作有限实数。

\begin{theorem}[Li 判据，引用]
RH 成立，当且仅当对每个正整数 $n$ 均有 $\lambda_n\ge0$。
这些系数也满足
\[
 \lambda_n=\left.
 \frac1{(n-1)!}\frac{\dd^n}{\dd s^n}
 \bigl(s^{n-1}\log\xi(s)\bigr)\right|_{s=1},
 \quad
 \xi(s)=\tfrac12s(s-1)\pi^{-s/2}\Gamma(s/2)\zeta(s).
\]
\end{theorem}
$s=1$ 附近 $\xi$ 不为零，可取局部解析对数；加入 $2\pi\ii$ 的常数，
因 $s^{n-1}$ 的 $n$ 阶导数为零，不改变定义。
零点式与导数式的一致性及“全部非负推出 RH”的反向定理在此引用，
上面直接证明的是 RH 推出非负性的方向。

\section{首个系数可以直接计算}
由导数式 $\lambda_1=\xi'(1)/\xi(1)$。
对 $\xi$ 取对数导数，在 $s=1$ 附近写
\[
 \frac{\xi'}{\xi}(s)=\frac1s+\frac1{s-1}
 -\frac12\log\pi+\frac12\frac{\Gamma'}{\Gamma}(s/2)
 +\frac{\zeta'}{\zeta}(s).
\]
利用 $\zeta(s)=(s-1)^{-1}+\gamma_E+O(s-1)$，
有 $\zeta'/\zeta(s)=-(s-1)^{-1}+\gamma_E+O(s-1)$。
又由 Gamma 倍元公式的对数导数，
$\Gamma'/\Gamma(1/2)=-\gamma_E-2\log2$。消去极点后，
\[
 \lambda_1=1+\frac{\gamma_E}{2}-\frac12\log(4\pi)>0.
\]
最后的不等号可用标准常数的可靠区间核查；近似值约为 $0.0230957$。
一个正系数只检查了一个必要条件，不是全部正整数上的结论。

\section{有限零点表留下怎样的尾项}
在 RH 下，对正虚部 $\gamma>0$，
\[
 \theta_\rho=2\arctan\frac1{2\gamma},\qquad
 0<\theta_\rho\le\frac1\gamma.
\]
用 $1-\cos v\le v^2/2$，得到共轭对贡献至多为 $n^2/\gamma^2$。
因此高度 $T$ 以上的尾和满足
\[
 0\le\sum_{\gamma>T}2(1-\cos(n\theta_\rho))
 \le n^2\sum_{\gamma>T}\gamma^{-2}
 =O\left(\frac{n^2\log(T+2)}T\right).
\]
计数估计给出的常数与 $n$ 无关，但若需要严格数值认证，仍须把它显式化。

关键限定是：这里使用了 RH，才能把所有尾项写成单位圆相位。
不能用这个条件性非负尾项去证明尚未假定的 RH。
只有临界线上的有限零点表，也不等于已经包含所有未知离线零点。

\section{等价改写的价值与边界}
判据的价值在于，把位置问题转到系数、导数及正性工具上。
若找到不依赖 RH 的方法证明对所有 $n$ 非负，就完成了所需的证明。
但对前若干个 $n$ 的数值验证，既需要控制各系数的计算误差，
也没有处理无穷多个其余指标。它是有限范围的检查，不是判据中全称量词的替代。

\section*{参考依据与后续核验}
X.-J.\ Li, \emph{The Positivity of a Sequence of Numbers and the Riemann Hypothesis},
Journal of Number Theory 65 (1997), 325--333，
\url{https://doi.org/10.1006/jnth.1997.2137}。
本书关于 $\xi$、零点计数及 RH 等价表述的结果可作为预备背景。
后续审阅重点为反向定理的引用范围、零点式的求和约定与尾项估计的条件性；
本篇不是该原论文的完整重证。
\end{document}
